\scanpage{30}

plongement normal) $=0$. Il faut donner un critère simple sur $\alpha$ pour qu'il en soit ainsi.

\uline{Question} Supposons que $h\in G$ soit tel que pour tte composante connexe $X_i$ de $X$, désignant par $\varphi_i:X_i\longrightarrow S$ le plongmt induit, $\pi_1(\deleted{\varphi}h\varphi_i)=\pi_1(h)\circ\pi_1(\varphi_i)\in\deleted{\unsure{\mathrm{Homext}}}\inserted{\mathrm{Hom\,ext}}(\pi_1(X_i),\pi_1(S))$ soit \unsure{égal|égale} à $\pi_1(\varphi_i)$, \deleted{Alors a-t-on} (condition évidemment nécessaire pour que \deleted{\unsure{$f$}} $h\circ\varphi$ soit \unsure{isoto|isot|isotope} à $\varphi$, \deleted{donc} i.e. $h\in H$) ; alors \unsure{à-t-on|a-t-on} $h\in H$ ?
% archive sheet number: 29
